%%%%%%%%%%%%%%%%%%%%%%%%%%%%%%%%%%%%%%%%%%%%%%%%%%%%%%%%%%%%%%%%%%%%%%%%%%%
% Handout und Uebungsblatt zu Lehreinheit 3
% Im Skript am Kapitelende, in handouts.tex als Druckvorlage.
% Der Kruskal-Text (Aufgabe 1) wird mit tools/kruskal_check.py geprueft.
% Wer ihn aendert, muss das Skript erneut laufen lassen!
%%%%%%%%%%%%%%%%%%%%%%%%%%%%%%%%%%%%%%%%%%%%%%%%%%%%%%%%%%%%%%%%%%%%%%%%%%%
\begin{handout}{Literaturrecherche}

\textbf{Worum es heute ging.} Alle haben eine eigene Startzahl gewählt und
sind nach derselben Zählregel durch einen Text gegangen, fast alle landeten
auf demselben Wort. Das ist Mathematik: Sobald zwei Wege einmal auf
dasselbe Wort treffen, laufen sie für immer gemeinsam weiter (Kruskal-Prinzip).
Eine systematische Recherche funktioniert ähnlich: Mit guten Suchbegriffen und
dem Weg über Literaturlisten landet man aus verschiedenen Richtungen bei den
zentralen Quellen. Die Kehrseite: Wer nur nach Bestätigung sucht, landet
zuverlässig bei der eigenen Vorannahme.

\textbf{Kernbegriffe}
\begin{description}
  \item[Systematische Recherche] geplante, dokumentierte und wiederholbare
    Suche mit festgelegten Begriffen, Datenbanken und Auswahlkriterien.
  \item[Boolesche Operatoren] \texttt{AND} (beide Begriffe), \texttt{OR}
    (einer genügt, für Synonyme), \texttt{NOT} (ausschließen, sparsam!).
  \item[Trunkierung und Phrase] \texttt{lern*} findet lernen, Lernende,
    Lernerfolg; \texttt{\symbol{34}magic trick\symbol{34}} sucht die genaue Wortfolge.
  \item[Schneeballsuche] rückwärts: Literaturliste eines Schlüsseltexts;
    vorwärts: Arbeiten, die ihn später zitieren \parencite{wohlin2014snowballing}.
  \item[Graue Literatur] nicht verlagsmäßig publiziert: Abschlussarbeiten,
    Preprints, Berichte. Oft aktuell, aber ohne Begutachtung.
  \item[Gegenrecherche] gezielte Suche nach Belegen, die der eigenen These
    widersprechen könnten (Congruence Bias, \textcite{wason1960failure}).
\end{description}

\textbf{Suchstring in drei Schritten}
\begin{center}
\small
\begin{tabular}{lll}
\toprule
\textbf{Konzept 1} & \textbf{Konzept 2} & \textbf{Konzept 3} \\
\midrule
\texttt{magic OR conjuring} & \texttt{wonder OR awe OR curiosity} & \texttt{learn* OR memory} \\
\bottomrule
\end{tabular}

\smallskip
\texttt{(K1) AND (K2) AND (K3)}: Synonyme mit OR, Konzepte mit AND
\end{center}

\begin{merkbox}
Eine Recherche ist erst dann wissenschaftlich, wenn jemand anderes sie mit
Ihrem Protokoll wiederholen könnte und wenn Sie auch nach dem gesucht haben,
was Ihnen widerspricht.
\end{merkbox}

\textbf{Zum Weiterlesen:} \textcite{doering2016forschungsmethoden}, Kapitel
zur Literaturrecherche; \textcite{gusenbauer2020academic} zur Frage, welche
Suchsysteme sich für systematische Recherchen eignen.
\end{handout}

\begin{uebungsblatt}{Recherchieren und dokumentieren}

\aufgabe{Der Zählweg}[in der Stunde, Einzelarbeit, 8 Min.]
Wählen Sie im Stillen eine Zahl zwischen 1 und 10. Zählen Sie vom Anfang des
Textes so viele Wörter ab; das ist Ihr Startwort. Zählen Sie seine Buchstaben
und gehen Sie genau so viele Wörter weiter. Wiederholen Sie das. Satzzeichen
zählen nicht, ä, ö, ü zählen als ein Buchstabe. Wenn der nächste Sprung über
das Textende hinausführen würde, ist Ihr aktuelles Wort Ihr
\textbf{Schlusswort}. Behalten Sie es für sich. \emph{(Text selbst verfasst,
fiktive Geschichte.)}

\medskip
\noindent\fbox{\parbox{\dimexpr\textwidth-2\fboxsep-2\fboxrule\relax}{%
\linespread{1.25}\selectfont\raggedright
\hyphenpenalty=10000 \exhyphenpenalty=10000 % keine Trennungen: jedes Wort ganz
% KRUSKAL-TEXT-BEGINN
Als Hanna Berger zum ersten Mal die alte Bibliothek ihrer Stadt betrat,
suchte sie nur ein einziges Buch. Ihre Tante hatte ihr erzählt, dass dort
ein Bericht über den Winter stehe, in dem der See vollständig zugefroren
war und die Menschen trocken bis zur anderen Seite gehen konnten. Hanna
wollte wissen, ob diese Geschichte stimmte. Sie fragte den Bibliothekar,
einen ruhigen Mann mit grauem Bart, wo sie anfangen solle. Er lächelte und
gab ihr keinen Titel, sondern einen Rat: \enquote{Fang irgendwo an, aber
schreib auf, wohin dich jedes Buch schickt.} Also begann sie bei einer
Chronik der Fischerei. Darin fand sie einen Hinweis auf alte
Kirchenbücher, die jede Taufe und jede Beerdigung festhielten. In den
Kirchenbüchern stand am Rand eine Notiz über ungewöhnliche Kälte, mit
einem Verweis auf die Wetterbeobachtungen eines Lehrers. Der Lehrer hatte
jeden Morgen die Temperatur gemessen und in kleine Hefte geschrieben.
Diese Hefte lagen im Keller, in einer Kiste mit der Aufschrift
\enquote{Verschiedenes}. Hanna arbeitete drei Wochen lang an einem langen
Tisch unter einem hohen Fenster. Sie las, verglich und notierte, und
manchmal fand sie Angaben, die einander widersprachen. Ein Kaufmann
schrieb, das Eis sei dünn gewesen, ein Pfarrer meinte, man habe darauf
getanzt. Hanna lernte, dass jeder Zeuge seine eigenen Gründe hatte, etwas
zu behaupten. Sie ordnete die Berichte nach Datum, nach Herkunft und nach
der Frage, wer was gesehen haben konnte. Am Ende führten fast alle Spuren,
gleich wo sie begonnen hatte, zu demselben Ort. Es war ein dünnes Heft
des Lehrers, in dem für einen Februar viele Nächte mit strengem Frost
standen und daneben ein Satz in sauberer Schrift: \enquote{Heute ging ich
mit meinen Schülern über das Eis bis zum anderen Ufer.} Hanna legte den
Stift hin. Sie wusste jetzt, dass die Erzählung ihrer Tante einen wahren
Kern hatte, und sie wusste noch etwas anderes: Wer lange und gründlich
sucht, landet oft bei derselben Quelle.
% KRUSKAL-TEXT-ENDE
}}

\medskip
\noindent Meine Startzahl: \rule{15mm}{0.3pt}\qquad
Mein Schlusswort: \rule{45mm}{0.3pt}\qquad
Zahl der Sprünge: \rule{15mm}{0.3pt}

\aufgabe{Suchbegriffe sammeln}[in der Stunde, Einzelarbeit, 10 Min.]
Zerlegen Sie Ihre Forschungsfrage aus LE~2 in zwei bis drei Konzepte. Sammeln
Sie zu jedem deutsche und englische Begriffe, Synonyme, Ober- und
Unterbegriffe. Markieren Sie, wo Trunkierung (\texttt{*}) sinnvoll ist.

\smallskip
\noindent\textbf{Meine Forschungsfrage:}
\schreiblinien{1}
\begin{center}
\renewcommand{\arraystretch}{2.0}
\begin{tabular}{|p{0.22\textwidth}|p{0.22\textwidth}|p{0.22\textwidth}|p{0.22\textwidth}|}
\hline
 & \textbf{Konzept 1} & \textbf{Konzept 2} & \textbf{Konzept 3} \\ \hline
Begriff deutsch & & & \\ \hline
Synonyme & & & \\ \hline
englisch & & & \\ \hline
Ober-/Unterbegriffe & & & \\ \hline
\end{tabular}
\end{center}

\aufgabe{Suchstring und Gegenrecherche}[in der Stunde, 5 Min.]
\begin{enumerate}
  \item Bauen Sie aus der Tabelle einen Suchstring: Synonyme mit \texttt{OR},
    Konzepte mit \texttt{AND}.
  \item Formulieren Sie einen zweiten Suchstring, der gezielt Arbeiten findet,
    die Ihrer erwarteten Antwort \emph{widersprechen} könnten (Begriffe wie
    \texttt{no effect}, \texttt{null result}, \texttt{failed replication},
    \texttt{critique}, oder die Gegenhypothese selbst).
\end{enumerate}
\noindent Suchstring 1:
\schreiblinien{2}
\noindent Suchstring 2 (Gegenrecherche):
\schreiblinien{2}

\aufgabe{Rechercheprotokoll}[in der Stunde beginnen, Hausaufgabe, 10\,\% der Portfolionote]
Führen Sie eine dokumentierte Recherche in mindestens zwei Datenbanken durch.
Eine Zeile pro Suchvorgang. Notieren Sie Suchbegriffe, Datenbanken,
Treffermengen und Auswahlkriterien so genau, dass eine andere Person Ihre
Suche wiederholen kann.

\begin{center}
\small
\renewcommand{\arraystretch}{2.4}
\begin{tabular}{|p{0.07\textwidth}|p{0.12\textwidth}|p{0.24\textwidth}|p{0.11\textwidth}|p{0.07\textwidth}|p{0.07\textwidth}|p{0.09\textwidth}|}
\hline
\textbf{Datum} & \textbf{Datenbank} & \textbf{Suchstring} & \textbf{Filter} &
\textbf{Tref\-fer} & \textbf{ge\-sich\-tet} & \textbf{aus\-ge\-wählt} \\ \hline
 & & & & & & \\ \hline
 & & & & & & \\ \hline
 & & & & & & \\ \hline
 & & & & & & \\ \hline
 & & & & & & \\ \hline
\end{tabular}
\end{center}
\noindent\textbf{Ein- und Ausschlusskriterien} (z.\,B. Sprache, Zeitraum,
Publikationstyp, Population):
\schreiblinien{2}
\noindent\textbf{Schneeball:} Welche zwei Schlüsseltexte haben Sie rückwärts
und vorwärts verfolgt? Was kam dabei neu hinzu?
\schreiblinien{2}

\aufgabe{Treffer sortieren}[in der Stunde, 10 Min.]
Ordnen Sie Ihre ersten fünf Treffer einer Stufe zu. Die Stufen sind eine
Arbeitshilfe für die erste Sichtung, kein Qualitätsurteil im Einzelfall; das
folgt in LE~4.

\begin{center}
\small
\begin{tabular}{cp{0.78\textwidth}}
\toprule
\textbf{Stufe} & \textbf{Publikationstyp} \\
\midrule
A & systematische Übersichtsarbeit, Metaanalyse, begutachteter Zeitschriftenartikel \\
B & Fachbuch, Handbuchkapitel, begutachteter Konferenzbeitrag \\
C & graue Literatur: Preprint, Abschlussarbeit, Bericht einer Behörde oder Stiftung \\
D & Fachpresse, Wissenschaftsjournalismus, Lehrbuchwissen ohne Belege \\
E & Blog, Video, Werbung, Quelle ohne erkennbare Autorschaft \\
\bottomrule
\end{tabular}
\end{center}

\aufgabe{Annotierte Bibliographie}[Hausaufgabe, Abgabe bis LE~4]
Kommentieren Sie fünf gefundene Quellen mit je zwei bis drei Sätzen zu
\emph{Kernaussage}, \emph{Methode} und \emph{Relevanz für Ihre Frage}.
Mindestens eine Quelle soll aus der Gegenrecherche stammen. Verwalten Sie die
Quellen in einem Literaturverwaltungsprogramm (Zotero, Citavi oder JabRef).

\smallskip
\noindent\textbf{Beispiel einer Annotation:}

\noindent\begin{minipage}{\textwidth}\small
Wason, P.~C. (1960). On the failure to eliminate hypotheses in a conceptual
task. \emph{Quarterly Journal of Experimental Psychology, 12}(3), 129--140.

\emph{Kernaussage:} Wer eine Regel erraten soll, prüft fast nur Beispiele, die
zur eigenen Vermutung passen, und bemerkt deshalb eine zu enge Hypothese
nicht. \emph{Methode:} Experiment mit dem Zahlentripel 2--4--6; nur sechs
von 29 Personen fanden die Regel ohne falschen Zwischenschluss.
\emph{Relevanz:} Begründet, warum meine Recherche auch nach widersprechenden
Befunden suchen muss.
\end{minipage}

\begin{lehrhinweis}[Erwartungshorizont]
Aufgabe~1: Schlusswort ist \textbf{\enquote{Quelle}} (letztes Wort des
Textes) für jede Startzahl von 1 bis 10; nach dem Prüfskript
\texttt{tools/kruskal\_check.py} laufen alle zehn Wege spätestens beim 47.~Wort
(\enquote{gehen}) zusammen, insgesamt sind es je nach Start 59 bis 66 Sprünge.
Wer woanders landet, hat sich verzählt; das gemeinsame Nachzählen ist eine
kleine Lektion in Reproduzierbarkeit. Aufgabe~2/3: mindestens zwei Konzepte,
englische Begriffe vorhanden, Synonyme mit OR innerhalb der Klammer, AND
zwischen Konzepten; Gegenrecherche enthält die Gegenhypothese oder
Begriffe für Nullbefunde und Kritik. Aufgabe~4: Protokoll ist
wiederholbar, wenn Datum, Datenbank, exakter String, Filter und Trefferzahl
dastehen; mindestens zwei Datenbanken, davon eine Fachdatenbank, Google
Scholar allenfalls ergänzend. Aufgabe~6: Annotationen trennen Kernaussage,
Methode und Relevanz; keine bloße Wiedergabe des Abstracts.
\end{lehrhinweis}
\end{uebungsblatt}
